\documentclass[12pt]{article}
\usepackage[utf8]{inputenc}
\usepackage[T1]{fontenc}
\usepackage{hyperref}
\title{Keerray-Woorroong Language Resources for K-12 Curriculum Development}
\author{Compiled by Claude Sonnet 5.5 (claude-sonnet-5-5) for Andr\'as Kornai}
\date{}
\begin{document}
\maketitle

\section{Introduction}
Keerray-Woorroong (also spelled Keerray Woorroong or Girai Wurrung) is an Aboriginal language of south-western Victoria, Australia, linked with the Gunditjmara people. Sources disagree on whether it is a separate language or a dialect of the wider Dhauwurd Wurrung (Gunditjmara) continuum.

\section{Primary Educational Resources}
The search reported that the Victorian Aboriginal Corporation for Languages (VACL) has a phrasebook, \emph{Yarnin Up: Keerray Woorroong}, with audio. It also reported VACL electronic story books on Gunditjmara people and languages. No URL for either was seen, so they are not in the JSON list.

\section{Linguistic and Reference Materials}
\begin{itemize}
\item AIATSIS AUSTLANG entry S25 for Keerray Woorroong. It lists alternative names and sources, including Krishna-Pillay (1996), \emph{Dictionary of Keerraywoorroong and related dialects}: \url{https://collection.aiatsis.gov.au/austlang/language/S25}
\item Wikipedia, Dhauwurd Wurrung language (the broader group): \url{https://en.wikipedia.org/wiki/Dhauwurd_Wurrung_language}
\item Museums Victoria, Victorian Aboriginal languages: \url{https://museumsvictoria.com.au/bunjilaka/about-us/first-peoples/victorian-aboriginal-languages/}
\item Warrnambool, Gunditjmara language page: \url{https://www.moyjil.com.au/gunditjmara-language}
\end{itemize}

\section{Religious and Community Resources}
The AUSTLANG entry points to the Gunditj Mirring Traditional Owners Aboriginal Corporation, and VACL is the main language organisation. No URLs for these were seen. The foreword to \emph{Indigenous Songs of Victoria} (ANU Press) is by Dr Vicki Couzens, a Keerray Wooroong Gunditjmara scholar, and discusses song and language revitalisation: \url{https://press-files.anu.edu.au/downloads/press/n14364/pdf/foreword.pdf}

\section{Digital and Audio Resources}
VACL's phrasebook audio is described above. No other audio, corpus or software resources were found.

\section{Curriculum Development Context}
The language is not known to be taught in schools, and revival work is led by the community and VACL. A curriculum would depend on community permission and on those organisations. The published dictionary and the VACL materials would be the starting points.
\end{document}